\section{Self-attention biểu diễn được tích chập}
\label{sec:ch10-attention-lam-duoc}

\index{thiên kiến quy nạp!attention biểu diễn được convolution}%
Tới đây thì câu chuyện có vẻ đã rõ: mạng \tn{tích chập}{convolution} có hai
ràng buộc, Transformer không có, nên Transformer thiếu một thứ. Chỗ này là chỗ
câu chuyện ấy hỏng.

Cordonnier, Loukas và Jaggi chứng minh rằng một tầng multi-head self-attention
\emph{biểu diễn được} một tầng tích chập bất kỳ, miễn đủ head. Định lý 1 của
bài~\autocite{cordonnier2020relationship}:

\begin{quote}
\enquote{A multi-head self-attention layer with \(N_h\) heads of dimension
\(D_h\), output dimension \(D_{out}\) and a relative positional encoding of
dimension \(D_p \geq 3\) can express any convolutional layer of kernel size
\(\sqrt{N_h} \times \sqrt{N_h}\) and \(\min(D_h, D_{out})\) output channels.}
\end{quote}

Đọc ngược lại thì điều kiện gọn: \textbf{cần \(N_h = K^2\) head cho một kernel
\(K \times K\)}. Một head cho mỗi độ lệch trong lân cận. Và positional encoding
phải là loại \emph{tương đối}, số chiều ít nhất 3.

Phép dựng đằng sau nó dễ hình dung hơn phát biểu. Chương~\ref{ch:transformer}
đã dựng \tn{ma trận điểm số}{score matrix} rồi lấy softmax để ra trọng số; nếu ép phân phối
attention của head thứ \(h\) dồn hết vào đúng một độ lệch cố định, thì head ấy
không làm gì khác ngoài đi lấy một pixel lân cận. Ghép \(K^2\) head như vậy lại
và trộn theo ma trận value, thì đó đúng là phép tích chập.

\begin{measured}{dựng tích chập bằng \(K^2\) head, so với \code{F.conv2d}}
\begin{minted}{text}
kernel  heads  max |heads - conv2d|  output scale  relative
3x3     9      7.629e-06             26.823        2.84e-07
5x5     25     1.717e-05             40.610        4.23e-07
7x7     49     2.670e-05             52.470        5.09e-07
\end{minted}

Đo bằng \code{experiments/ch10\_equivariance.py} tại tag \code{ch10}.
\end{measured}

Sai lệch tuyệt đối lớn lên theo kernel, và cột cuối cho thấy nó lớn nhanh hơn
biên độ đầu ra một chút: 2.84 lên 5.09 phần mười triệu khi số head đi từ 9 lên
49. Đó đúng là dáng của sai số cộng dồn float32, vì phép dựng cộng \(K^2\) số
hạng thay vì một, và không phải dáng của một khoảng cách cấu trúc - một khoảng
cách cấu trúc không nằm ở mức \(10^{-7}\) tương đối và không đi theo số số hạng
được cộng.

\textbf{Và tôi phải nói rõ phép đo ấy đo cái gì.} Nó đo \emph{sức biểu diễn},
không đo việc học. Các head ở đây tôi đặt bằng tay, ở giới hạn softmax đã
\tn{bão hòa}{saturation} hẳn. Dựng một tầng self-attention thật mà softmax bão
hòa được thì phải đẩy hệ số của \tn{mã hóa vị trí}{positional encoding} bậc hai
lên đủ lớn, và khi ấy con số báo cáo việc
softmax bão hòa tới đâu chứ không báo cáo định lý đúng hay sai. Bảng trên là
điểm cuối của một phép dựng, không phải kết quả của một phép huấn luyện.

Chính bài cũng cẩn thận đúng chỗ ấy, khi nói về trường hợp một chiều:

\begin{quote}
\enquote{Since we have not tested empirically if the preceding construction
matches the behavior of 1D self-attention in practice, we cannot claim that it
actually learns to convolve an input sequence - only that it has the capacity
to do so.}
\end{quote}

\subsection*{Cái định lý ấy đổi của chương này}

\index{thiên kiến quy nạp!khác biệt không nằm ở sức biểu diễn}%
Nếu self-attention biểu diễn được tích chập, thì khác biệt giữa hai kiến trúc
\emph{không} nằm ở chỗ cái này làm được điều cái kia không làm được.

Nó nằm ở chỗ cái nào là mặc định.

Một tầng tích chập \emph{bắt buộc} phải cục bộ và dùng chung trọng số: không có
cấu hình tham số nào của nó phá được hai ràng buộc ấy, vì hai ràng buộc nằm
trong cách tầng được viết ra chứ không nằm trong giá trị của trọng số. Một tầng
self-attention có thể trở thành một tầng tích chập, nhưng nó phải \emph{học}
để trở thành, từ dữ liệu, cùng lúc với việc học mọi thứ khác.

Nói theo cách của bài ViT: hai ràng buộc kia được \enquote{baked into each
layer}. Nướng sẵn vào. Cái nướng sẵn thì không tốn dữ liệu; cái phải học thì
tốn.

Đó là toàn bộ luận điểm của chương này, và mục~\ref{sec:ch10-anh-that} là chỗ
tôi đặt một con số lên chữ \enquote{tốn}.

Bài Cordonnier còn đo thêm một chuyện nói đúng hướng ấy: các tầng đầu của một
mạng chỉ-attention huấn luyện trên ảnh \emph{tự} học ra kiểu chú ý theo lưới
quanh mỗi pixel, tức tự đi tới chỗ mà tích chập bắt đầu từ đó. Bài kết luận:

\begin{quote}
\enquote{Our results seem to suggest that localized convolution is the right
inductive bias for the first few layers of an image classifying network.}
\end{quote}

Một kiến trúc không có \tn{thiên kiến quy nạp}{inductive bias} về ảnh, cho đủ
dữ liệu, học lại được chính cái thiên kiến ấy. Mục sau là bốn lần người ta thử
xoay quanh chỗ đó: ba lần đặt sẵn một phần của nó vào thay vì bắt mô hình học,
và một lần không đặt gì cả rồi trả bằng tính toán.
